\documentclass[11pt,letterpaper]{gradescopeexam}
%% NOTE: This file contains Python code which is evaluated
%% by a program `pythontex`, like this:
%%     shell% pythontex Midterm1
%% This is only relevant if \printanswers is in effect
%% The Pythontex package can be found here:
%%     https://github.com/gpoore/pythontex
\noprintanswers
%\printanswers
%\printsolutionkey

\usepackage{mdframed}
\usepackage[right,modulo]{lineno}
\usepackage{pythontex}

\usepackage{amsmath}
\usepackage{bbm}
\usepackage{amsthm}
\usepackage{amsfonts}
\usepackage{times}
\usepackage{graphicx}
\usepackage{multicol}
\usepackage{amsthm}
\usepackage{datetime}
\usepackage{hyperref}
\usepackage{comment}
\usepackage{listings}
\usepackage{xcolor}
\usepackage{units}
\usepackage{enumitem}
% \usepackage{MnSymbol}
\usepackage{pgfplots}
\usepackage{hyperref}


\renewcommand\P{\mathbb{P}}
\newcommand\R{\mathbb{R}}
\newcommand\Q{\mathbb{Q}}
\newcommand\E{\mathbb{E}}
\newcommand\N{\mathbb{N}}
\newcommand\F{\mathbb{F}}
\newcommand\one{\mathbbm{1}}
\newcommand\V{\mathrm{var}}
\newcommand\sgn{\mathrm{sgn}}
\renewcommand\vec[1]{\mathbf{#1}}

\renewcommand\c[1]{\mathcal{#1}}
\newcommand\im[1]{\mathrm{Im}(#1)}
\DeclareMathOperator\Cov{\mathrm{Cov}}

\newtheorem*{lemma}{Lemma}
\newtheorem*{theorem}{Theorem}

\usepackage{geometry}
\usepackage{listings} % Include listings package

% Setup the listings package for Python syntax highlighting
\lstset{language=Python,
  basicstyle=\ttfamily\small,
  keywordstyle=\color{blue},
  commentstyle=\color{olive},
  stringstyle=\color{red},
  showstringspaces=false,
  identifierstyle=\color{black}}

% 
% Define style for the Maxima language listings
% 
\lstdefinelanguage{Maxima}{
  keywords={addrow,addcol,zeromatrix,ident,augcoefmatrix,ratsubst,diff,ev,tex,%
    with_stdout,nouns,express,depends,load,submatrix,div,grad,curl,%
    rootscontract,solve,part,assume,sqrt,integrate,abs,inf,exp},
  morekeywords=[2]{for,while,if,then,thru,do,define,subst,ev,makelist,sum,nusum,int},  
  morekeywords=[3]{plot2d,plot3d,wxplot2d,wxplot3d},
  morekeywords=[4]{alpha,beta,gamma,delta,eta,lambda},
  sensitive=true,
  comment=[n][\itshape]{/*}{*/},
  commentstyle=\color{olive},
  literate=
  *{+}{{{\color{purple}+}}}1
  {-}{{{\color{purple}-}}}1
  {:}{{{\color{purple}:}}}2
  {:=}{{{\color{purple}:=}}}1
  {*}{{{\color{purple}*}}}1
  {^}{{{\color{purple}\^{}}}}1    
}


\usepackage{algorithm}
\usepackage{algorithmic}


\newcommand{\continuationbox}[0]{%
  \newpage
  \noindent%
  \textbf{\sc Continuation of Problem:} \fbox{\vphantom{$\int$}\phantom{Q00.00}}
  \begin{mdframed}[shadow=true,
    linecolor=blue,
    linewidth=1pt,
    shadowcolor=black!40,
    ]%
    \vspace{0.9\textheight}
  \end{mdframed}}


\pgfplotsset{compat=1.18}

\begin{document}

% ----------------------------------------------------------------
\LeftBoxMargin{0.0\linewidth}
\RightBoxMargin{0.0\linewidth}
\AssignmentTitle{Midterm 1 (Max.~score = \numpoints\ points)}
\CourseName{Math 589B, Spring 2024}
\ifprintanswers
\FirstName{Marek}
\LastName{Rychlik}
\StudentId{31415926}
\fi
% ----------------------------------------------------------------
\makeheader
\vspace{0.1in}

\begin{questions}
  \begin{question}[10]
    Given is the following Newton's divided difference table of the function \(f(x)\):
    \[
    \begin{array}{c|ccccc}
      x_i & f[x_i] & f[x_i,x_{i+1}] & f[x_i,x_{i+1},x_{i+2}] 
      & f[x_i,x_{i+1},x_{i+2},x_{i+3}] 
      & f[x_i,x_{i+1},x_{i+2},x_{i+3},x_{i+4}] \\[6pt]
      \hline
      0      & 0       & 0.7956 & 0.1312 & 0.2379 & 0.1792 \\
      0.25   & 0.1989  & 0.8612 & 0.3096 & 0.4171 &        \\
      0.50   & 0.4142  & 1.016  & 0.6224 &        &        \\
      0.75   & 0.6682  & 1.3272 &        &        &        \\
      1.00   & 1       &        &        &        &        \\
    \end{array}
    \]
    Use the 4th‑degree Newton polynomial constructed from the table to
    approximate \( f(0.1) \).  Make sure to give the answer with the
    correct number of significant digits, assuming that the values in
    the table are correctly rounded to \(4\) significant digits, i.e.,
    have relative error \(\le 0.5\times 10^{-4}\) (\(1/2\) ulp guarantee).
\begin{pycode}
I= 0 \
   + 0.7956 * (0.1 - 0) \
   + 0.1312 * (0.1 - 0) * (0.1-0.25) \
   + 0.2379 * (0.1-0) * (0.1 - 0.25)  * (0.1 - 0.5) \
   + 0.1792 * (0.1-0) * (0.1 - 0.25)  * (0.1 - 0.5) * (0.1 -0.75)
n = 10
def r(x):
    return round(x,n)
\end{pycode}
 
    \begin{solutionorbox}[4in]
      \[
      \begin{aligned}
        f(0.1) \approx \\
        & f[0] \\
        & + f[0,0.25]\,(0.1-0) \\
        & + f[0,0.25,0.50]\,(0.1-0)\,(0.1-0.25) \\
        & + f[0,0.25,0.50,0.75]\,(0.1-0)\,(0.1-0.25)\,(0.1-0.50) \\
        & + f[0,0.25,0.50,0.75,1]\,(0.1-0)\,(0.1-0.25)\,(0.1-0.50)\,(0.1-0.75) \\
        &= 0 
          + \underbrace{0.7956\,(0.1)}_{\py{r(0.7956*0.1)}}
          + \underbrace{0.1312\,(0.1)(-0.15)}{\py{r(0.1312*(0.1))}}
          + \underbrace{0.2379\,(0.1)(-0.15)(-0.40)}_{\py{r(0.2379*(0.1)*(-0.15)*(-0.40))}} \\
        & + \underbrace{0.1792\,(0.1)(-0.15)(-0.40)(-0.65)}_{\py{r(0.1792*(0.1)*(-0.15)*(-0.40)*(-0.65))}} \\
        &= \py{I} \approx\py{round(I,5)}.
      \end{aligned}
      \]
      The result of the computation is exact, but every divided
      difference has 4 digits after the decimal point, and the
      \(1/2\)-ulp guarantee assumes that all divided difference
      approximations are within \(0.5\cdot 10^{-4}\) of the exact
      values. Since the firs value in the table is multiplied by
      \(0.1\), which is known exactly. Therefore
      \(0.7956\cdot 0.1 = 0.07956 \) has \(6\) accurate digits, i.e.,
      \(5\) accurate digits after the decimal point.

      The other quantities in the sum have negligible impact
      on accuracy of the sum because they are by at least an order of magnitude
      smaller. Hence, we should round the answer to \(5\) digits
      after the decimal point.

      \TheAnswer{\(\py{round(I,5)}\), rounded to \textbf{five digits} after the decimal point}
    \end{solutionorbox}
  \end{question}    
  
  \begin{question}
    In numerical integration, the Midpoint Rule (for a single
    subinterval) approximates the definite integral of a function $f$
    over $[a,b]$ by evaluating $f$ at the midpoint $m = \frac{a+b}{2}$:
    \[
    \int_a^b f(x)\,dx 
    \;\approx\; 
    (b - a)\, f\Bigl(\frac{a + b}{2}\Bigr).
    \]

    \begin{parts}
      \part[10]
      Write down (in your own words) the statement of the Midpoint Rule
      for approximating an integral of a continuous function $f$ over
      the interval $[a,b]$.
      \begin{prooforbox}[5.7in]
        The formula results from approximating the area by a rectangle
        of height \(f(m)\) and base \((b-a)\). Drawing in the book.
      \end{prooforbox}
      \part[10]
      Using ideas from how the error was derived for the Trapezoidal
      Rule (which we discussed in class), derive an error expression of
      the form
      \[
      \text{Error} 
      \;=\; 
      -\,\frac{(b-a)^3}{24}\,f''(\xi)
      \]
      for some $\xi \in (a,b)$.  (Hint: consider the second derivative
      of $f$, apply an appropriate form of the Remainder Theorem or
      Taylor expansion, and carefully compare the exact integral to the
      midpoint approximation.)
      \begin{solutionorbox}[5.5in]
        We use Taylor expansion about the middle point \(m\):
        \[ f(x) = f(m) + f'(m)(x-m) + \frac{f''(\xi_x)}{2}(x-m)^2. \]
        Integrating over \(x\) over \([a,b]\) gives
        \[
        \begin{aligned}
 \int f(x)\,dx &=
          f(m)(b-a) + f'(m)\overbrace{\int_a^b(x-m)\,dx}^{0} + \int_a^b\frac{f''(\xi_x)}{2}(x-m)^2\,dx\\
          &= f(m)(b-a) + \frac{f''(\xi)}{2} \int_a^b(x-m)^2\,dx.
        \end{aligned}
        \]
        In the last step we applied the Mean Value Theorem for integrals.
        Finally,
        \[ 
        \begin{aligned}
          \int_a^b(x-m)^2\,dx
          &= \frac{1}{3}(x-m)^3\bigg|_{x=a}^b\\
          &= \frac{1}{3}\left((b-m)^3-(a-m)^3\right)\\
          &= \frac{1}{3} \left(
            \left(\frac{b-a}{2}\right)^3-
            \left(\frac{a-b}{2}\right)^3
            \right)\\
          &= \frac{1}{3} \left(
            \left(\frac{b-a}{2}\right)^3+
            \left(\frac{b-a}{2}\right)^3
            \right)\\
          &= \frac{1}{3} \left(
            \left(\frac{b-a}{2}\right)^3+
            \left(\frac{b-a}{2}\right)^3
            \right)\\
          &= \frac{2}{3} \left(\frac{b-a}{2}\right)^3\\
          &= \frac{1}{12} (b-a)^3.
        \end{aligned}
        \]
        Combining with the previous, we obtain:

        \TheAnswer{\(\text{Error} = \frac{f''(\xi)}{24}(b-a)^3\)}

        NOTE: In particular, when \(f\) is convex, the error is
        positive, which means that the Midpoint Rule underestimates
        the integral (the opposite of the Trapezoid Rule).
      \end{solutionorbox}
      \part[10] 
      Give a detailed proof that shows why there must exist such a
      point $\xi$ in $(a,b)$ satisfying the error bound. (You may use
      the Mean Value Theorem for integrals, Rolle s Theorem or other
      standard theorems covered in class and/or calculus.)
      \begin{prooforbox}[6.5in]
        The standard ``proof'' is given in the derivation of the previous part.
        However, the proof is known to have a gap, because \(\xi_x\)
        may not depend continuously on \(x\).
\begin{comment}
        The correct proof results if we use the integral form of the remainder
        in the Taylor series:
        \[ f(x) = f(m) + f'(m)(x-m) + R_1(x) \]
        where
        \[ \begin{aligned}
          R_1(x) &= \int_m^x\frac{f''(t)}{2!}(x-t)^2\,dt\\
                 &=\frac{1}{2}\int_0^1 f''(m+u(x-m))(x-m-u(x-m))^2(x-m)\,du\\
                 &=\frac{1}{2}(x-m) \int_0^1 f''(m+u(x-m))((1-u)(x-m))^2\,du\\
                 &=\frac{1}{2}(x-m)^3 \int_0^1 f''(m+u(x-m))(1-u)^2\,du.
        \end{aligned}
        \]


        Therefore, after integrating,
        \[
        \begin{aligned}
          \text{Error}
          &= \int_a^b R_1(x)\,dx\\
          & = \int_a^b\left(\frac{1}{2}(x-m)^3 \int_0^1 f''(m+u(x-m))(1-u)^2\,du\right),dx\\
          & = \frac{1}{2}\iint_{[a,b]\times[0,1]}(x-m)^3 f''(m+u(x-m))(1-u)^2\,du\,dx\\
          & = \frac{1}{2}\int_0^1 f''(m+u(x-m))(1-u)^2(x-m)^3 \,dx\,du.\\
        \end{aligned}
        \]

  \]
  Where 
  \[
  \begin{aligned}
    \int_m^b \frac{f''(t)}{2} \left(\int_t^b(x-t)^2\,dx\right)\,dt\\
    &= \frac{f''(\xi)}{2} \int_m^b\int_t^b(x-t)^2\,dx\,dt.
  \end{aligned}
  \]
  We find by exact calculation that
  \[ \int_m^b\int_t^b(x-t)^2\,dx\,dt = \frac{(b-a)^3}{24}.\]
\end{comment}
        
      \end{prooforbox}
      % \part[10]
      % Explain how this error formula depends on the interval length
      % $(b-a)$ and on the curvature of $f$. Compare it to the error
      % constant in the Trapezoidal Rule.
      % \begin{prooforbox}[3in]
      % \end{prooforbox}
    \end{parts}
  \end{question}

  \begin{question}
    Consider the Legendre polynomial of degree 3, $P_{3}(x)$. By
    definition, Legendre polynomials can be obtained from the
    \emph{Rodrigues formula}:
    \[
    P_{n}(x) \;=\; \frac{1}{2^{n}\,n!}\;\frac{d^{n}}{dx^{n}} \bigl[(x^{2}-1)^{n}\bigr].
    \]
    \begin{parts}
      \part[10] Use the Rodrigues formula with $n=3$ to derive an explicit
      expression for $P_{3}(x)$.
      \begin{solutionorbox}[5.8in]
        We start with the Rodrigues formula:
        \[
        P_{3}(x)
        \;=\; \frac{1}{2^3\,3!}\,\frac{d^3}{dx^3}\Bigl[(x^2-1)^3\Bigr]
        \;=\; \frac{1}{8\,6}\,\frac{d^3}{dx^3}\bigl[(x^2-1)^3\bigr].
        \]
        First, expand $(x^2-1)^3 = (x^2-1)^2\,(x^2-1) = (x^4 - 2x^2 + 1)\,(x^2 - 1)$. One may also expand directly:
        \[
        (x^2-1)^3 
        \;=\; x^6 - 3x^4 + 3x^2 - 1.
        \]
        Then,
        \[
        \frac{d^3}{dx^3}\bigl(x^6 - 3x^4 + 3x^2 - 1\bigr)
        \;=\;
        \frac{d^3}{dx^3}\bigl(x^6\bigr)
        \;-\;3\,\frac{d^3}{dx^3}\bigl(x^4\bigr)
        \;+\;3\,\frac{d^3}{dx^3}\bigl(x^2\bigr)
        \;-\;\frac{d^3}{dx^3}(1).
        \]
        Compute term by term:
        \[
        \frac{d^3}{dx^3}(x^6) = 6\cdot5\cdot4\,x^3 = 120\,x^3, 
        \quad
        \frac{d^3}{dx^3}(x^4) = 4\cdot3\cdot2\,x = 24\,x,
        \]
        \[
        \frac{d^3}{dx^3}(x^2) = 2\cdot1\cdot0\, x^0 = 0,
        \quad
        \frac{d^3}{dx^3}(1)=0.
        \]
        Hence,
        \[
        \frac{d^3}{dx^3}\bigl[(x^2-1)^3\bigr]
        =120\,x^3 \;-\;3\,(24\,x)
        =120\,x^3 \;-\;72\,x
        =24\,(5x^3 - 3x).
        \]
        Therefore,
        \[
        P_{3}(x)
        \;=\;
        \frac{1}{48}\;\bigl[\,24\,(5x^3 - 3x)\bigr]
        \;=\;
        \frac{5x^3 - 3x}{2}.
        \]
        \TheAnswer{
          \[ P_3(x) = \frac{5x^3 - 3x}{2} \]
        }

      \end{solutionorbox}

      \part[10] Find the roots of $P_{3}(x)$. These roots will be the
      \emph{nodes} $x_{1} < x_{2} < x_{3}$ for the 3-point
      Gauss--Legendre quadrature rule on the interval $[-1,1]$.
      \begin{solutionorbox}[6.5in]
        Set $P_{3}(x)=0$. Then
        \[
        \frac{5x^3 - 3x}{2} \;=\;0
        \quad\Longrightarrow\quad
        x\,(5x^2 - 3)\;=\;0.
        \]
        This gives the solutions
        \[
        x=0,
        \quad
        x = \pm\,\sqrt{\frac{3}{5}}.
        \]
        Thus, the three nodes are
        \[
        x_1 \;=\; -\,\sqrt{\frac{3}{5}}, 
        \quad
        x_2 \;=\; 0,
        \quad
        x_3 \;=\; \sqrt{\frac{3}{5}}.
        \]
      \end{solutionorbox}

      \part[10] The standard formula for the \emph{weights} $w_{i}$ in
      Gauss--Legendre quadrature is
      \[
      w_{i} 
      \;=\; \frac{2}{\bigl(1 - x_{i}^{2}\bigr)\,\bigl[P_{n}'(x_{i})\bigr]^{2}},
      \]
      where $P_{n}'(x)$ is the derivative of the Legendre polynomial of degree $n$. 
      Compute the weights $w_1, w_2, w_3$ for $n=3$, using the roots $x_i$ from part (b).
      \begin{solutionorbox}[6in]
      We use the formula
      \[
      w_i \;=\; \frac{2}{\bigl(1 - x_i^2\bigr)\,\bigl[P_{3}'(x_i)\bigr]^2}.
      \]
      First, differentiate $P_{3}(x)$:
      \[
      P_{3}'(x)
      \;=\;
      \frac{d}{dx}\Bigl(\tfrac{5x^3 - 3x}{2}\Bigr)
      \;=\;
      \frac{15x^2 - 3}{2}
      \;=\;
      \frac{3\bigl(5x^2 - 1\bigr)}{2}.
      \]

      \underline{\emph{Weight at $x_2=0$.}}
      \[
      P_{3}'(0) 
      = \frac{3\bigl(5\cdot0^2 - 1\bigr)}{2}
      = -\frac{3}{2}.
      \]
      Hence,
      \[
      w_2
      = \frac{2}{\bigl(1 - 0^2\bigr)\,\bigl(-\tfrac{3}{2}\bigr)^2}
      = \frac{2}{1 \cdot \tfrac{9}{4}}
      = \frac{2}{\tfrac{9}{4}}
      = \frac{8}{9}.
      \]

      \underline{\emph{Weights at $x_1=-\sqrt{\frac{3}{5}}$ and $x_3=\sqrt{\frac{3}{5}}$.}}

      By symmetry, these weights are equal. Let $x_+ = \sqrt{\frac{3}{5}}$. Then
      \[
      1 - x_+^2
      = 1 - \frac{3}{5}
      = \frac{2}{5},
      \quad
      P_{3}'(x_+)
      = \frac{3}{2}\,\bigl(5\cdot\tfrac{3}{5}-1\bigr)
      = \frac{3}{2}\,(3-1)
      = \frac{3}{2}\cdot 2
      = 3.
      \]
      Thus
      \[
      \bigl[P_{3}'(x_+)\bigr]^2
      = 9,
      \quad
      w_1 = w_3
      = \frac{2}{\left(\tfrac{2}{5}\right)\,9}
      = \frac{2}{\tfrac{18}{5}}
      = \frac{2 \cdot 5}{18}
      = \frac{5}{9}.
      \]
      \end{solutionorbox}
      
      \part[10] List your final answers for the \emph{nodes} and
      \emph{weights} $\{(x_1,w_1),\,(x_2,w_2),\,(x_3,w_3)\}$ in
      ascending order of the $x_i$. Hence, your completed 3-point
      Gauss--Legendre rule will have the form
      \[
      \int_{-1}^{1} f(x)\,dx 
      \;\approx\;
      \sum_{i=1}^{3} w_i\,f\bigl(x_i\bigr).
      \]
      \begin{solutionorbox}[6in]
    \TheAnswer{
      \[
      \begin{aligned}
        &x_1 = -\sqrt{\tfrac{3}{5}},  &w_1 =\tfrac{5}{9},\\
        &x_2 = 0, &w_2 = \tfrac{8}{9},\\
        &x_3 = \sqrt{\tfrac{3}{5}}, &w_3 =\tfrac{5}{9}.
      \end{aligned}
      \]}
    \end{solutionorbox}
    \end{parts}
  \end{question}


  \ifprintanswers
  \else
    \continuationbox
    \continuationbox
  \fi




\end{questions}
\end{document}
